\documentclass[exercice,a4]{thedocument}%
\usepackage[french]{babel}%
\usepackage{pstricks}%
\usepackage{tabularx}%
\usepackage{array}%
\usepackage{eurosym}%
\usepackage{enumitem}%
\startdatakeys
    \source{22GENMATAS1}
    \BO{2026}
    \cycle{4}
    \competencesDuSocle{CA,RA}
    \objectifsApprentissage{C4B1bc}
    \motsCles{volume,prix,conversion}
\enddatakeys
\begin{document}%

\medskip

On souhaite construire un carré potager en utilisant des planches en bois et en suivant le montage ci-dessous.

Le carré potager souhaité n'a pas de fond et il a la forme d'un pavé droit de base carrée et de hauteur $30$~cm.

\begin{center}
\begin{tabularx}{\linewidth}{|c|X|}\hline
\textbf{Vue de dessus} &\textbf{Plan et indications pour le montage}

Prévoir dans chaque angle une équerre à visser avec 8 vis pour assembler les 4 planches formant l'angle.\\
\psset{unit=1cm}
\begin{pspicture}(4.8,5)
\psframe(0,0.6)(4.4,4.9)
\psline(0.2,0.8)(0.2,4.7)\psline(4.2,0.8)(4.2,4.7)
\psline(0,4.7)(4.4,4.7)
\psline(0,0.8)(4.2,0.8) \psline(4.2,0.6)(4.2,4.7)
\psline[linestyle=dashed]{<->}(0,0.4)(4.2,0.4)\uput[d](2.1,0.4){1,20~m}
\psline[linestyle=dashed]{<->}(4.4,0.4)(4.2,0.4)\uput[d](4.3,0.4){2 cm}
\end{pspicture}&\psset{unit=0.9cm}
\begin{pspicture}(9.5,4.5)
\psframe(0,0.8)(5.2,2.1)\psline(0,1.45)(5.3,1.45)(7.9,3.75)\psline(7.1,3.75)(2.6,3.75)(0.8,2.1)
\pspolygon(5.3,0.8)(7.9,3.1)(7.9,4.4)(5.3,2.1)
\psline(7.9,4.4)(2.6,4.4)(0,2.1)
\psline(5.2,2.1)(7.7,4.3)(2.6,4.3)(0.15,2.1)
\psline(2.6,4.3)(2.6,3.1)
\psline(6.3,3.1)(2.6,3.1)(1.6,2.1)
\psline(5.2,2.1)(5.3,2.1)\psline(5.2,0.8)(5.3,0.8)
\rput(0.7,4.2){2 planches}
\psline{->}(0.7,4)(2.2,3.6)
\psline{->}(0.7,4)(2,2.8)
\psline[linestyle=dashed]{<->}(0,0.6)(5.2,0.6)\uput[d](2.6,0.6){1,20~m}
\psline[linestyle=dashed]{<->}(5.5,0.7)(8.1,3.1)\uput[dr](6.8,1.9){1,20~m}
\psline[linestyle=dashed]{<->}(8.1,3.1)(8.1,4.4)\uput[r](8.1,3.75){30~cm}
\rput(5.4,0.4){2 cm}\rput(8.4,2.7){2 cm}
\end{pspicture}\\ \hline
\end{tabularx}
\begin{tabularx}{\linewidth}{|*{4}{>{\centering \arraybackslash}X|}}\hline
\multicolumn{4}{|c|}{\textbf{Prix}}\\\hline
Équerre à 8 trous &Planche en bois&Vis&\small Sac de terre végétale\\
&{\footnotesize 250 cm $\times$ 15 cm $\times$ 2 cm}&Lot de 100&40 L\\
2,90~\euro{} la pièce&5,60~\euro{} la pièce&5,70~\euro{} le lot&6,90~\euro{} le sac\\ \hline
\end{tabularx}
\end{center}

\smallskip

\begin{enumerate}
\item À l'achat, les planches en bois mesurent $2,50$~m de longueur.
	\begin{enumerate}
		\item Combien de planches devra-t-on acheter ?
		\item Déterminer le budget nécessaire (hors coût de la terre) pour réaliser ce carré potager.
	\end{enumerate}
\end{enumerate}

On remplit le carré potager de terre végétale au minimum jusqu'aux deux tiers de sa hauteur.

On dispose la terre afin qu'elle forme un pavé droit dont la longueur du côté de la base carrée est de $118$~cm.

\begin{enumerate}[resume]
\item Sept sacs de terre végétale seront-ils suffisants pour compléter au minimum le carré potager ?

On rappelle que : 1~L = 1dm$^3$.
\end{enumerate}

\showthesource
\end{document}
